\documentclass[11pt]{article}
\usepackage[margin=1in]{geometry}
\usepackage{amsmath,amssymb}
\usepackage{enumitem}
\usepackage{bookmark}
\usepackage{xurl}
\hypersetup{
  pdftitle={PSET 6: Discrete random variables},
  hidelinks}

\title{PSET 6: Discrete random variables}
\author{}
\date{\vspace{-2.5em}}

% Problems are numbered 1-7 across the document.
\newcounter{problem}
\newcommand{\problem}[1]{\stepcounter{problem}\section*{Problem \theproblem: #1}}

\begin{document}
\maketitle

% Shared assignment questions (header block).
% Include at the top of any problem set with:  % Shared assignment questions (header block).
% Include at the top of any problem set with:  % Shared assignment questions (header block).
% Include at the top of any problem set with:  \input{assignment-qs}
\section*{Assignment Qs}

\begin{itemize}
\item Name
\item How long did this problem set take you (round to nearest hour)? \quad
  0-1 hour \ 2-3 hours \ 4-5 hours \ 5+ hours
\item How difficult was this problem set? \quad
  very easy \ 1 \ 2 \ 3 \ 4 \ 5 \ very challenging
\end{itemize}

\section*{Assignment Qs}

\begin{itemize}
\item Name
\item How long did this problem set take you (round to nearest hour)? \quad
  0-1 hour \ 2-3 hours \ 4-5 hours \ 5+ hours
\item How difficult was this problem set? \quad
  very easy \ 1 \ 2 \ 3 \ 4 \ 5 \ very challenging
\end{itemize}

\section*{Assignment Qs}

\begin{itemize}
\item Name
\item How long did this problem set take you (round to nearest hour)? \quad
  0-1 hour \ 2-3 hours \ 4-5 hours \ 5+ hours
\item How difficult was this problem set? \quad
  very easy \ 1 \ 2 \ 3 \ 4 \ 5 \ very challenging
\end{itemize}


\problem{Calculate probabilities in a sample space $S$}

Events $A$ and $B$ are contained within a sample space $S$. Given that
$\Pr(A) = 0.65$, $\Pr(B) = 0.3$, and $\Pr(A \cap B) = 0.1$,
find:\footnote{Inspired by Grimmer HW8.4}

\begin{enumerate}[label=\alph*.]
\item $\Pr(A \cup B)$
\item $\Pr(A \cap B^c)$
\item $\Pr[(A \cap B^c) \cup (B \cap A^c)]$
\end{enumerate}

\problem{Random variables}

Provide the following explanations in your own words:

\begin{enumerate}[label=\alph*.]
\item What is a random variable?
\item What is the difference between upper case and lower case $x$? How
  (if at all) do they matter?
\item If I'm thinking about something like $p_X(x_0)$, what do all the
  parts/pieces mean?
\end{enumerate}

\problem{Survey says}

A survey has 54\% respondents 50 or older and 46\% respondents under 50.
Within the survey, on a particular question, 9.5\% of the 50-plus
population agrees strongly while 2.7\% of under-50 respondents agree
strongly.

\begin{enumerate}[label=\alph*.]
\item What is the probability someone selected at random is 50 or older?
\item The selected individual strongly agrees with the survey question.
  Now what is the likelihood that person is 50 or older? Explain your
  reasoning and show all your work.
\item Are the two answers above the same or different? Explain.
\item (for fun, no points) What is the survey question?
\end{enumerate}

\problem{PMF vs.\ CMF}

Consider the following function: $f(x) = \frac{1}{6}$. Find the pmf and
cmf of the function and provide them in a table below.

\problem{Getting a traffic ticket}

You drive to work 5 days a week for a full year (50 weeks), and with
probability $p = 0.04$ you get a traffic ticket on any given day,
independent of other days. Let $X$ be the total number of tickets you
get in the year.\footnote{Inspired by BT 2.41}

\begin{enumerate}[label=\alph*.]
\item What is the probability that the number of tickets you get is
  exactly equal to the expected value of $X$?
\item Calculate approximately the probability in (a) using a Poisson
  approximation.
\end{enumerate}

\problem{Obtaining requests for information}

$X$ is a discrete random variable. It takes the value of the number of
days required for a governmental agency to respond to a request for
information. $X$ is distributed according to the following
PMF:\footnote{Inspired by Grimmer HW10.2}
\[
f(x) = e^{-6} \dfrac{6^x}{x!} \quad \mbox{for} \ x \in \{0, 1, 2, \ldots\}
\]

\begin{enumerate}[label=\alph*.]
\item Given this information, what is the probability of a response from
  the agency in 5 days or less?
\item What is the probability the agency response takes more than 10 but
  less than 13 days?
\item What is the probability the agency response takes more than 5
  days?
\item Suppose using $X$ you generate a new variable, \textbf{Responsive}.
  \textbf{Responsive} equals 1 if an agency responds in 5 days or less
  and 0 otherwise. What is the expected value of \textbf{Responsive}?
\item What is the variance of \textbf{Responsive}?
\end{enumerate}

\problem{Modeling electoral outcomes}

Suppose we've developed a model predicting the outcome of the upcoming
midterm elections in a state with 4 Congressional districts. In each
district there are two candidates, a Republican and a Democrat. We have
reason to believe the following PMF describes the distribution of
potential election results, where $K \in \{0, 1, 2, 3, 4\}$ is the number
of seats won by Republican candidates in the upcoming election.
\[
\Pr(K = k \mid \theta) = \binom{4}{k} \theta^{k} (1-\theta)^{4-k}
\]

Based on polling information, we think the appropriate value for $\theta$
is 0.423.\footnote{Inspired by Grimmer HW10.3. Data from
  \url{https://www.realclearpolling.com/polls/state-of-the-union/generic-congressional-vote}}

\begin{enumerate}[label=\alph*.]
\item What's the expected number of seats Republicans will win in the
  upcoming election?
\item Given this PMF, what's the probability that no Republican
  legislators win in the upcoming election?
\item What's the probability that Republican legislators win a majority
  of the seats in this state?
\item A prominent political pundit declares they are certain that
  Republicans will win a majority of seats in the next election and
  offers the following bet. If Republicans win a majority of the seats,
  we must pay the pundit \$15.00. If Republicans fail to win a majority
  of the seats, we will win \$20.00. Based on our model, should we take
  this bet? \textbf{Hint: think of the betting outcomes as a random
  variable. Find the expected value of this random variable.}
\item Suppose we are offered a second bet with a more complicated
  structure. In this case we'll receive \$100 if the Republicans win a
  majority, \$50 if neither party wins a majority, and we'll have to pay
  \$200 if the Democrats win a majority. Should we take this bet?
\end{enumerate}

% Shared AI and Resources statement.
% Include in any problem set with:  % Shared AI and Resources statement.
% Include in any problem set with:  % Shared AI and Resources statement.
% Include in any problem set with:  \input{ai-statement}
\section*{AI and Resources statement}

Please list (in detail) all resources you used for this assignment. If
you worked with people, list them here as well. It is not enough to say
that you used a resource for help, you need to be specific on the link
and \emph{how} it was helpful. W/R/T gen AI tools (including GPT,
Claude, etc.) you cannot use them to do work on your behalf---you
cannot put in any of the questions, etc. You can ask for help on logic
/ sample problems. If you do use GPT or other AI tools, you need to
provide a link to your chat transcript. Any suspected academic
integrity violations will be immediately reported to the Dean of
Students.

\section*{AI and Resources statement}

Please list (in detail) all resources you used for this assignment. If
you worked with people, list them here as well. It is not enough to say
that you used a resource for help, you need to be specific on the link
and \emph{how} it was helpful. W/R/T gen AI tools (including GPT,
Claude, etc.) you cannot use them to do work on your behalf---you
cannot put in any of the questions, etc. You can ask for help on logic
/ sample problems. If you do use GPT or other AI tools, you need to
provide a link to your chat transcript. Any suspected academic
integrity violations will be immediately reported to the Dean of
Students.

\section*{AI and Resources statement}

Please list (in detail) all resources you used for this assignment. If
you worked with people, list them here as well. It is not enough to say
that you used a resource for help, you need to be specific on the link
and \emph{how} it was helpful. W/R/T gen AI tools (including GPT,
Claude, etc.) you cannot use them to do work on your behalf---you
cannot put in any of the questions, etc. You can ask for help on logic
/ sample problems. If you do use GPT or other AI tools, you need to
provide a link to your chat transcript. Any suspected academic
integrity violations will be immediately reported to the Dean of
Students.


\end{document}
